% !TEX program = xelatex
% !TEX root = eigenfaces_algorithm.tex
% 编译（推荐，自动处理目录所需的多次编译）：
%   latexmk -xelatex eigenfaces_algorithm.tex
% 或手动：
%   xelatex eigenfaces_algorithm.tex
%   xelatex eigenfaces_algorithm.tex
% 依赖：TeX Live / MiKTeX，需 ctex、amsmath、booktabs、hyperref、geometry。
% Overleaf：编译器选 XeLaTeX，主文件设为本文件。
\documentclass[12pt]{ctexart}

\usepackage[a4paper,margin=2.4cm]{geometry}
\usepackage{amsmath,amssymb}
\usepackage{booktabs}
\usepackage{xcolor}
\usepackage[colorlinks=true,linkcolor=blue!60!black,urlcolor=blue!60!black]{hyperref}

\newcommand{\mT}{^{\mathsf{T}}}
\newcommand{\R}{\mathbb{R}}
\newcommand{\mat}[1]{\mathbf{#1}}
\newcommand{\norm}[1]{\left\lVert #1 \right\rVert}
\newcommand{\inner}[2]{\left\langle #1,\,#2 \right\rangle}

% “直觉解释”环境：给初学者的白话段落
\newenvironment{intuition}
{\par\medskip\noindent\hspace{-0.5em}{\bfseries 【直觉解释】}\quad}
{\par\medskip}

% “为什么这样做”环境
\newenvironment{why}
{\par\medskip\noindent\hspace{-0.5em}{\bfseries 【为什么要这一步】}\quad}
{\par\medskip}


\title{\textbf{特征脸（Eigenfaces）算法详解}\\[4pt]
\large 幂迭代、PCA 重建与识别——从零开始的完整推导}
\author{配合《高级工程数学》第 7 题\ 研究任务使用}
\date{2026 年 11 月}

\begin{document}
\maketitle

\begin{abstract}
本文档不假设读者会 PCA 或特征分解的任何“套路”，从“机器如何认人”这个问题出发，
一步一步讲清楚特征脸算法的每一环节：为什么要中心化、为什么可以在 $280\times 280$ 的小矩阵上做、
幂迭代为什么收敛、收缩为什么必须再正交化、重建和识别到底在算什么。
每个公式都先给白话直觉，再给严格推导。
\end{abstract}

\tableofcontents
\bigskip

% !TEX root = ../eigenfaces_algorithm.tex
%=====================================================
\section{背景：我们想让机器“认人”，但它只会算数}

\subsection{问题是什么}

我们手里有 \textbf{Olivetti Faces} 人脸库：40 个人，每人 10 张灰度照片，
每张照片是 $64\times 64$ 个像素，每个像素是一个 $0\sim 1$ 之间的亮度数字。
任务有两个：

\begin{enumerate}
  \item \textbf{重建}：用很少的“信息”把一张人脸大致拼回来；
  \item \textbf{识别}：拿来一张没见过的人脸，判断它是 40 个人里的哪一位。
\end{enumerate}

\subsection{第一步：把照片变成向量}

计算机眼里没有“脸”，只有数字。把 $64\times 64$ 的像素网格按行拉直，
就得到一个 $4096$ 维的列向量：
\[
\mathbf{a}=\begin{pmatrix}a_1\\ a_2\\ \vdots\\ a_{4096}\end{pmatrix}\in\R^{4096}.
\]
于是“一个人”在这个数据库里的身份，变成了 $\R^{4096}$ 空间里的一个点。
40 个人 $\times$ 10 张照片，一共 400 个点。

\begin{intuition}
把每个人脸想成 $\R^{4096}$ 这个超高维空间里的一个“萤火虫”。
40 个人就是 40 簇萤火虫，每簇 10 只。识别问题就是：飞来一只新萤火虫，
看它离哪一簇最近。整个特征脸算法，本质上就是在想办法
“把灯光调暗，只留下萤火虫，让簇与簇分得更开”。
\end{intuition}

\subsection{为什么不能直接比较像素？}

最天真的做法：新照片 $\mathbf{a}$ 和库里的照片 $\mathbf{b}$ 算欧氏距离
$\norm{\mathbf{a}-\mathbf{b}}$，最近的就是答案。问题是：

\begin{itemize}
  \item 同一个人，光照从左边换成右边，两张照片的像素距离可能比\textbf{不同两个人}的距离还大；
  \item 4096 个维度里，绝大多数像素（背景、额头、脸颊）几乎是“死”的——不携带身份信息；
  \item 距离被这些无信息维度平均稀释，噪声淹没信号。
\end{itemize}

\begin{intuition}
辨认朋友靠的不是 4096 个像素点，而是几个关键特征：脸型轮廓、五官位置、发际线。
也就是说，\textbf{虽然照片活在 4096 维里，但“人的差异”只活在很少的几个维度里}。
特征脸算法的全部目标，就是自动找出这几个维度。
\end{intuition}

% !TEX root = ../eigenfaces_algorithm.tex
%=====================================================
\section{预备知识：三件趁手的工具}

后面的推导只依赖三件事，这里一次性复习。

\subsection{内积：衡量“相似程度”的尺子}

两个向量的内积
\[
\inner{\mathbf{x}}{\mathbf{y}}=\mathbf{x}\mT\mathbf{y}
=\sum_i x_i y_i=\norm{\mathbf{x}}\,\norm{\mathbf{y}}\cos\theta ,
\]
其中 $\theta$ 是两向量的夹角。三个重要特例：
\begin{itemize}
  \item $\inner{\mathbf{x}}{\mathbf{y}}=0$：两个方向互相\textbf{垂直}（正交），彼此毫无共同成分；
  \item $\norm{\mathbf{x}}=1$（单位向量）时，$\inner{\mathbf{x}}{\mathbf{y}}$ 就是 $\mathbf{y}$ 在 $\mathbf{x}$ 方向上的\textbf{投影长度}，也就是“$\mathbf{y}$ 里含有多少 $\mathbf{x}$ 的成分”；
  \item $\norm{\mathbf{x}}=\norm{\mathbf{y}}=1$ 时，内积就是夹角余弦：$1$ 表示同向，$0$ 表示垂直。
\end{itemize}

\subsection{矩阵乘向量：拉伸机器}

$\mat{A}\mathbf{x}$ 的效果是把 $\mathbf{x}$ 旋转 + 拉伸。一般情况下方向会变；
但存在一些特殊方向，乘完之后\textbf{方向不变，只被拉长了 $\lambda$ 倍}：
\[
\mat{A}\mathbf{v}=\lambda \mathbf{v}.
\]
$\mathbf{v}$ 叫特征向量，$\lambda$ 叫特征值。$\lambda$ 的大小表示
“$\mat{A}$ 在这个方向上的拉伸力度”。

\subsection{对称矩阵的谱定理（本课程的母定理）}

若 $\mat{A}=\mat{A}\mT$ 实对称，则
\[
\text{所有特征值都是实数，且可以取到一组互相\textbf{正交}的单位特征向量。}
\]
即存在正交矩阵 $\mat{V}=[\mathbf{v}_1,\dots,\mathbf{v}_d]$（$\mat{V}\mT\mat{V}=\mat{I}$）使
\[
\mat{A}=\mat{V}\mat{\Lambda}\mat{V}\mT,\qquad
\mat{\Lambda}=\operatorname{diag}(\lambda_1\ge\lambda_2\ge\dots\ge\lambda_d).
\]

\begin{intuition}
对称矩阵没有“旋转”这种花活，它只会在若干个互相垂直的方向上各自拉伸。
这组互相垂直的拉伸方向（特征向量）配上拉伸倍数（特征值），
就是矩阵的完整“体检报告”。正交性还送我们一份大礼：
$\mat{V}^{-1}=\mat{V}\mT$，逆变换不需要解方程，转置一下就行。
\end{intuition}

% !TEX root = ../eigenfaces_algorithm.tex
%=====================================================
\section{数据组织：矩阵、均值脸与中心化}

\subsection{训练集与划分}

把 400 张照片按人分层划分：每人 10 张里固定取 7 张进训练集、3 张进测试集。
训练集共 $n=280$ 张，堆成一个数据矩阵：
\[
\mat{A}=\begin{pmatrix}
\rule{4pt}{9pt}\ \mathbf{a}_1^{\mathsf{T}}\ \rule{4pt}{9pt}\\
\rule{4pt}{9pt}\ \mathbf{a}_2^{\mathsf{T}}\ \rule{4pt}{9pt}\\
\vdots\\
\rule{4pt}{9pt}\ \mathbf{a}_{280}^{\mathsf{T}}\ \rule{4pt}{9pt}
\end{pmatrix}\in\R^{280\times 4096},
\qquad n=280,\quad d=4096 .
\]
每一\textbf{行}是一张脸（拉直后的向量转置），每一\textbf{列}是“某个像素位置在 280 张脸里的取值”。

\subsection{均值脸与中心化}

对 280 张训练脸逐像素取平均，得到\textbf{均值脸} $\bar{\mathbf{a}}=\frac1n\sum_{j}\mathbf{a}_j$。
把每张脸减去它：
\[
\mat{X}=\mat{A}-\mathbf{1}\,\bar{\mathbf{a}}\mT ,
\]
即 $\mat{X}$ 的每一行都是“这张脸减去平均长相”的\textbf{残差脸}。
（$\mathbf{1}$ 是 280 维全 1 向量，作用是把同一个均值脸复制 280 份逐行减掉。）

\begin{why}
如果不减均值会怎样？所有像素都非负、整体亮度贡献了一个巨大的“公共成分”，
数据云的重心远离原点。此时方差最大的方向几乎就是“平均亮度方向”，
它对区分谁是谁毫无帮助。中心化把云的重心挪到原点，
剩下的方差才反映“人与人、脸与脸的差异”。
\end{why}

\begin{why}
均值只用训练集算。测试集若用自己的均值，等于让模型提前“看了”测试数据（信息泄漏），
而且测试脸的坐标系会和训练出的特征脸对不上。这是本题第一坑。
\end{why}

\subsection{一个决定一切的不等式}

中心化后每行和为定值（都减了同一个均值），280 行中只有 $n-1=279$ 行独立，故
\[
\operatorname{rank}(\mat{X})\le n-1=279\ll 4096 .
\]
\begin{intuition}
280 张残差脸不管看起来多不一样，它们全体都躺在 4096 维空间里一个
\textbf{至多 279 维的薄板}上。也就是说，4096 个维度里至少有 3817 个是“空房间”。
我们不需要在 4096 维里找东西，只需要描述这块薄板——这就是后面一切“小矩阵技巧”的底气。
\end{intuition}

% !TEX root = ../eigenfaces_algorithm.tex
%=====================================================
\section{PCA：找“差异最大”的那几个方向}

\subsection{目标函数：方差 = 信息}

我们想找单位向量 $\mathbf{u}$（$\norm{\mathbf{u}}=1$），
把所有残差脸投影到它上面后，投影值尽可能“散”。投影值是 $t_j=\inner{\mathbf{u}}{\mathbf{x}_j}=\mathbf{x}_j\mT\mathbf{u}$，
其平方和（中心化后均值自动为零，这就是方差的两倍尺度，差常数因子不影响结果）为
\[
\sum_{j=1}^{n} t_j^2
=\sum_j \mathbf{u}\mT\mathbf{x}_j\,\mathbf{x}_j\mT\mathbf{u}
=\mathbf{u}\mT\Big(\underbrace{\mat{X}\mT\mat{X}}_{d\times d}\Big)\mathbf{u}.
\]

\begin{intuition}
为什么“方差大 = 信息多”？想象一个班里所有人身高完全相同——身高这个维度方差为零，
告诉你身高也猜不出是谁。反过来，一个维度上方差很大，说明这个维度上人与人有明显差异，
\textbf{差异正是识别所需要的信息}。PCA 的哲学就是：把方差大的维度挑出来，
把方差小（近乎人人都一样）的维度扔掉。
\end{intuition}

\subsection{化成特征值问题（三行推导）}

带约束的最大化问题：
\[
\max_{\mathbf{u}}\ \mathbf{u}\mT\mat{X}\mT\mat{X}\,\mathbf{u}
\qquad\text{s.t.}\qquad \mathbf{u}\mT\mathbf{u}=1 .
\]
用 Lagrange 乘子法，令
\[
L(\mathbf{u},\lambda)=\mathbf{u}\mT\mat{X}\mT\mat{X}\mathbf{u}-\lambda(\mathbf{u}\mT\mathbf{u}-1),
\]
对 $\mathbf{u}$ 求导置零（利用 $\nabla_{\mathbf{u}}\,\mathbf{u}\mT\mat{C}\mathbf{u}=2\mat{C}\mathbf{u}$，$\mat{C}$ 对称）：
\[
2\,\mat{X}\mT\mat{X}\,\mathbf{u}-2\lambda\mathbf{u}=\mathbf{0}
\quad\Longrightarrow\quad
\boxed{\ \mat{X}\mT\mat{X}\,\mathbf{u}=\lambda\mathbf{u}\ }.
\]
结论漂亮得过分：
\begin{itemize}
  \item 让方差最大的方向，\textbf{恰好是} $\mat{X}\mT\mat{X}$ 的特征向量；
  \item 把它代回目标函数：$\mathbf{u}\mT(\lambda\mathbf{u})=\lambda$，
        \textbf{最大方差值恰好就是那个特征值}。
\end{itemize}
所以“按特征值从大到小排”，前 $k$ 个特征向量就是信息量最大的 $k$ 个方向。
这些 4096 维的特征向量 reshaped 回 $64\times 64$ 去看，就是传说中的
\textbf{特征脸（eigenfaces）}：前几张常像“整体明暗/光照方向”，后面几张像五官与纹理的细碎模式。

\begin{intuition}
把 $\mat{X}\mT\mat{X}$ 想成“像素与像素之间的相关性报告”。
它是对称阵，由谱定理知道它有一套互相垂直的拉伸方向。
PCA 做的事，说穿了就是\textbf{问这份报告：“你力气最大的拉伸方向是哪几个？”}
\end{intuition}

% !TEX root = ../eigenfaces_algorithm.tex
%=====================================================
\section{关键困难与 Gram 技巧：在 280 那一侧算}

\subsection{困难：矩阵太大}

$\mat{C}=\frac{1}{n-1}\mat{X}\mT\mat{X}\in\R^{4096\times4096}$。
对它做完整特征分解约需 $O(d^3)\approx 7\times10^{10}$ 次运算，
光把它存下来就要 $4096^2\times 8\text{B}\approx 134$ MB。
而每个特征向量才用一次——太奢侈了。

\subsection{Gram 技巧：三个等式解决问题}

改看另一个矩阵：
\[
\mat{G}=\mat{X}\mat{X}\mT\in\R^{280\times 280}.
\]
它小得多。它的元素有直接含义：$G_{ij}=\inner{\mathbf{x}_i}{\mathbf{x}_j}$
是第 $i$ 张残差脸和第 $j$ 张残差脸的\textbf{相似度}——所以 $\mat{G}$ 叫 Gram 矩阵（相似度矩阵）。

\textbf{断言：}若 $\mat{G}\mathbf{u}=\lambda\mathbf{u}$ 且 $\lambda>0$，
则 $\mathbf{v}=\mat{X}\mT\mathbf{u}$ 是 $\mat{X}\mT\mat{X}$ 的特征向量，特征值同为 $\lambda$。

\textbf{证明（两行）：}
\[
\mat{X}\mT\mat{X}\,(\mat{X}\mT\mathbf{u})
=\mat{X}\mT\,(\mat{X}\mat{X}\mT\,\mathbf{u})
=\mat{X}\mT(\lambda\mathbf{u})
=\lambda\,(\mat{X}\mT\mathbf{u}).\qquad\blacksquare
\]

\textbf{归一化：}$\mat{X}\mT\mathbf{u}$ 不是单位向量，需要除以它的长度：
\[
\norm{\mat{X}\mT\mathbf{u}}^2
=\mathbf{u}\mT\mat{X}\mat{X}\mT\mathbf{u}
=\lambda\,\mathbf{u}\mT\mathbf{u}=\lambda
\quad\Longrightarrow\quad
\boxed{\ \mathbf{v}=\frac{\mat{X}\mT\mathbf{u}}{\sqrt{\lambda}}\ },\qquad \norm{\mathbf{v}}=1 .
\]

\subsection{两套谱的关系与“空房间”}

$\mat{X}\mT\mat{X}$ 与 $\mat{X}\mat{X}\mT$ 的非零特征值\textbf{完全相同}：
$\mat{X}$ 的秩 $\le 279$，所以两边都至多 279 个非零特征值。
协方差版本只差常数：$\mat{C}$ 的特征值为 $\lambda_i/(n-1)$。
此外若 $\mathbf{u}_i\perp\mathbf{u}_j$（Gram 侧），则
$\inner{\mathbf{v}_i}{\mathbf{v}_j}=\tfrac{1}{\sqrt{\lambda_i\lambda_j}}\mathbf{u}_i\mT\mat{X}\mat{X}\mT\mathbf{u}_j
=\tfrac{\lambda_j}{\sqrt{\lambda_i\lambda_j}}\inner{\mathbf{u}_i}{\mathbf{u}_j}=0$：
\textbf{特征脸彼此也是正交的}——正交性在小矩阵那侧算完就自动继承到大矩阵这侧。

\begin{intuition}
一个直观比喻：$\mat{X}$ 是一座桥，左岸住着 280 个“人”（$\mathbf{u}$ 们），
右岸住着 4096 个“像素方向”（$\mathbf{v}$ 们）。
桥上真正通行的“变异模式”只有 $\le 279$ 条。
在左岸（人这一侧）把这条条模式找清楚，桥的对称性会自动告诉我们右岸的对应者是谁：
$\mathbf{v}=\mat{X}\mT\mathbf{u}/\sqrt{\lambda}$。
这就是“$n\ll d$ 时先算 $n\times n$ 那一侧”的全部秘密。
\end{intuition}

\subsection{与 SVD 的对应（一张对照表）}

瘦奇异值分解 $\mat{X}=\mat{U}\mat{\Sigma}\mat{V}\mT$ 中
$\sigma_i=\sqrt{\lambda_i}$，于是：
\[
\text{Gram 侧：}\ \mat{G}=\mat{U}\mat{\Lambda}\mat{U}\mT,\qquad
\text{特征脸：}\ \mat{V}=\mat{X}\mT\mat{U}\mat{\Lambda}^{-1/2}.
\]
一句话总结本题的算法本质：\textbf{“特征脸 = 只用 $n$ 侧信息求出 $\mat{X}$ 的右奇异向量。”}

% !TEX root = ../eigenfaces_algorithm.tex
%=====================================================
\section{幂迭代：不分解矩阵，只“反复拉伸”}

\subsection{核心直觉：反复乘 $\mat{G}$，大方向自己会冒头}

取任意向量 $\mathbf{u}_0$，反复做“乘 $\mat{G}$ 再归一化”：
\[
\mathbf{u}_1=\frac{\mat{G}\mathbf{u}_0}{\norm{\mat{G}\mathbf{u}_0}},\qquad
\mathbf{u}_2=\frac{\mat{G}\mathbf{u}_1}{\norm{\mat{G}\mathbf{u}_1}},\qquad\dots
\]
为什么会收敛到最大特征向量？把 $\mathbf{u}_0$ 在（正交的）特征基上展开：
\[
\mathbf{u}_0=c_1\mathbf{u}_1^{(\star)}+c_2\mathbf{u}_2^{(\star)}+\cdots ,
\]
（记号：$\mathbf{u}_i^{(\star)}$ 是真实特征向量，区别于迭代中的 $\mathbf{u}_k$。）
乘 $k$ 次后，每个分量被各自的特征值放大 $\lambda_i^k$ 倍：
\[
\mat{G}^k\mathbf{u}_0
=\lambda_1^k\Big(c_1\mathbf{u}_1^{(\star)}
+\sum_{i\ge 2}c_i\Big(\frac{\lambda_i}{\lambda_1}\Big)^{\!k}\mathbf{u}_i^{(\star)}\Big).
\]
因为 $0\le\lambda_i/\lambda_1<1$，除最大方向外的所有成分按公比 $\lambda_i/\lambda_1$
\textbf{几何衰减}，最后只剩 $\mathbf{u}_1^{(\star)}$。
归一化只负责防止数值溢出（$\lambda_1^k$ 会爆表或归零），不影响方向。

\subsection{一个能手算的 $2\times 2$ 例子}

取 $\mat{G}=\begin{pmatrix}6.5&2.5\\2.5&6.5\end{pmatrix}$，
其特征值为 $\lambda_1=9$（方向 $(1,1)/\sqrt2$）、$\lambda_2=4$（方向 $(1,-1)/\sqrt2$）。
从 $\mathbf{u}_0=(1,0)$ 出发，两个方向的初始成分相等，
每迭代一步，“错误方向”与“正确方向”的成分比恰好乘上 $\lambda_2/\lambda_1=4/9$。
令 $\theta_k$ 为第 $k$ 步迭代向量与真实主方向的夹角，则 $\sin\theta_k\approx(4/9)^k$：

\begin{center}
\begin{tabular}{ccccccc}
\toprule
$k$ & 1 & 2 & 3 & 4 & 5 & 10\\
\midrule
$(4/9)^k$ & 0.444 & 0.198 & 0.088 & 0.039 & 0.017 & $3.0\times10^{-4}$\\
$\theta_k$（度）& $26.4^\circ$ & $11.4^\circ$ & $5.0^\circ$ & $2.2^\circ$ & $1.0^\circ$ & $0.017^\circ$\\
\bottomrule
\end{tabular}
\end{center}

\begin{intuition}
每乘一次 $\mat{G}$，正确方向被放大 9 倍、错误方向只放大 4 倍，
两者的差距像滚雪球一样拉大——5 步之后错误方向就只剩 $1.7\%$。
幂迭代没有任何魔法，它就是\textbf{让矩阵自己一次次投票}，票多的方向赢。
\end{intuition}

\subsection{收敛速度的定量结论}

由上面的展开立得（$\theta_k$ 为迭代向量与 $\mathbf{u}_1^{(\star)}$ 的夹角）：
\[
\sin\theta_k \lesssim \frac{c_2}{c_1}\Big(\frac{\lambda_2}{\lambda_1}\Big)^{k}
\qquad\Longrightarrow\qquad
\text{达到精度 }\varepsilon\text{ 需要 }k\approx\frac{\ln(\varepsilon\, c_1/c_2)}{\ln(\lambda_2/\lambda_1)} .
\]
两个直接推论：
\begin{itemize}
  \item \textbf{谱间隙 $\lambda_1-\lambda_2$ 越大收敛越快}；$\lambda_2/\lambda_1\to 1$ 时迭代次数爆炸——
  这正是本题“前 50 个用幂迭代、50 以后交给 \texttt{eigh}”的定量原因（尾部特征值密密麻麻挤在一起，比值趋近 1）。
  \item 初始向量几乎垂直于 $\mathbf{u}_1^{(\star)}$（$c_1\approx 0$）时会“慢启动”。
  随机高维向量恰好垂直的概率为零，但 $c_1$ 偏小会让前几步看起来停滞——
  这解释了换随机种子时迭代次数的小幅波动。
\end{itemize}

\subsection{算法步骤（完整版）}

求 $\mat{G}\in\R^{n\times n}$ 的最大特征对：
\begin{enumerate}
  \item 初始化：$\mathbf{u}_0\sim\mathcal{N}(\mathbf{0},\mat{I}_n)$，归一化 $\mathbf{u}_0\leftarrow\mathbf{u}_0/\norm{\mathbf{u}_0}$；固定随机种子保证可复现。
  \item 循环：
  \[
  \mathbf{w}\leftarrow \mat{G}\,\mathbf{u}_k,\qquad
  \mathbf{u}_{k+1}\leftarrow \frac{\mathbf{w}}{\norm{\mathbf{w}}_2}.
  \]
  \item 每步用 \textbf{Rayleigh 商}估计特征值：
  \[
  \hat\lambda_k=\mathbf{u}_k\mT\mat{G}\,\mathbf{u}_k .
  \]
  \item 停机准则：相对残差
  \[
  \frac{\norm{\mat{G}\mathbf{u}_k-\hat\lambda_k\mathbf{u}_k}_2}{\hat\lambda_k}<\varepsilon
  \quad(\text{建议 }10^{-8}\sim10^{-10}).
  \]
\end{enumerate}

\begin{why}
为什么特征值用 Rayleigh 商而不是 $\norm{\mat{G}\mathbf{u}_k}$？
Rayleigh 商在特征向量附近是\textbf{二阶精确}的：若方向误差为 $\delta$，
则 $|\hat\lambda-\lambda_1|=O(\delta^2)$——把 $\mathbf{u}_k=c_1\mathbf{u}_1^{(\star)}+\delta\,\mathbf{z}$
（$\mathbf{z}\perp\mathbf{u}_1^{(\star)}$）代入即可验证。
而 $\norm{\mat{G}\mathbf{u}_k}$ 只有一阶精度。所以 Rayleigh 商收敛“快人一步”，
代价是它太乐观：$\hat\lambda$ 看起来不动了，方向可能还差 $\sqrt\varepsilon$。
因此停机必须看\textbf{残差}而不是看 $\Delta\hat\lambda$。
\end{why}

% !TEX root = ../eigenfaces_algorithm.tex
%=====================================================
\section{收缩（Hotelling Deflation）：抹掉第一名，再选第二名}

\subsection{问题与做法}

幂迭代只给最大特征对。要第二个，就把已经找到的方向\textbf{从矩阵里抹掉}：
\[
\mat{G}_1=\mat{G}-\lambda_1\mathbf{u}_1\mathbf{u}_1\mT,
\qquad\text{然后对 }\mat{G}_1\text{ 重新跑幂迭代，得到 }(\lambda_2,\mathbf{u}_2).
\]
依此类推：$\mat{G}_2=\mat{G}_1-\lambda_2\mathbf{u}_2\mathbf{u}_2\mT$，……，直到攒够 $k$ 对。

\subsection{为什么这样抹是干净的（两行验证）}

新矩阵 $\mat{G}_1$ 仍对称。验证它在两个方向上的行为：
\begin{itemize}
  \item 沿 $\mathbf{u}_1$：$\mat{G}_1\mathbf{u}_1=\mat{G}\mathbf{u}_1-\lambda_1\mathbf{u}_1(\mathbf{u}_1\mT\mathbf{u}_1)=\lambda_1\mathbf{u}_1-\lambda_1\mathbf{u}_1=\mathbf{0}$——该方向特征值变为 $0$；
  \item 沿任何 $\mathbf{w}\perp\mathbf{u}_1$：$\mat{G}_1\mathbf{w}=\mat{G}\mathbf{w}-\lambda_1\mathbf{u}_1(\mathbf{u}_1\mT\mathbf{w})=\mat{G}\mathbf{w}$——完全不动。
\end{itemize}
所以 $\mat{G}_1$ 的谱就是“把 $\lambda_1$ 换成 $0$，其余不变”，
它的最大特征对恰是 $(\lambda_2,\mathbf{u}_2)$。\textbf{收缩 = 精准阉割一个方向，其余谱原封不动。}

\subsection{必须加的保险：逐步再正交化}

\begin{intuition}
理论上抹干净了，实际中抹不干净。计算机是浮点机器：
$\lambda_1\mathbf{u}_1\mathbf{u}_1\mT$ 这一项带着 $10^{-16}$ 量级的舍入误差残留。
下一次做 matvec 时这点残留被乘上 $\mat{G}$，放大 $\lambda_1/\lambda_i$ 倍——
对于尾部小特征值，这个放大倍数是天文数字。若干轮之后，
迭代向量会被“吸”回 $\mathbf{u}_1$，求出的“$\mathbf{u}_2$”其实还是第一名。
\end{intuition}

保险做法（正交化版幂迭代）：每次 matvec 之后，先把结果对已求出的所有方向投影剔除，再归一化：
\[
\mathbf{w}\leftarrow\mat{G}\mathbf{u}_k;\qquad
\mathbf{w}\leftarrow\mathbf{w}-\sum_{j\le i-1}\inner{\mathbf{u}_j}{\mathbf{w}}\,\mathbf{u}_j
\quad(\text{做两遍更稳});\qquad
\mathbf{u}_k\leftarrow\mathbf{w}/\norm{\mathbf{w}} .
\]
这与隐式收缩 $(\mat{I}-\mat{Q}\mat{Q}\mT)\mat{G}$ 数学等价，但数值上稳得多。
另外每轮的初始向量要重新随机取（上一轮的 $\mathbf{u}$ 正好躺在被抹掉的方向上，复用它没有意义）。

\subsection{误差累积：为什么只算前 50 个}

两个互补的机理，写报告时值得完整叙述：

\begin{enumerate}
  \item \textbf{除以间隙}：特征向量对扰动的敏感度 $\approx\ \norm{\text{扰动}}/(\lambda_i-\lambda_{i+1})$。
  头部间隙巨大（$90/95\%$ 方差集中在前几十个），尾部间隙趋于 0——
  同样的舍入扰动，在尾部被放大成巨大的方向误差。
  \item \textbf{噪声地板}：$i$ 次收缩后累积的算子误差 $\sim i\cdot\varepsilon_{\text{mach}}\cdot\lambda_1$。
  当 $\lambda_i$ 降到这个量级以下，幂迭代会在噪声里打转，返回的方向是噪声方向而非真实特征方向。
\end{enumerate}

于是实验里会观察到一个自洽的现象：$|\lambda_{i+1}/\lambda_i|\to 1$（间隙消失）
与第 $i$ 对的迭代次数暴涨、方向误差抬升\textbf{同步发生}。
“$k\le 50$ 用幂迭代 + 收缩，$k>50$ 用 \texttt{eigh}”不是偷懒，
而是算法在尾部已经到达可分辨极限。
\texttt{eigh} 用的是完全不同的正交三角化技术（$O(n^3)$ 对 280 阶矩阵瞬间完成），
它一次给出全部特征对且精度均匀，正好补上尾部的洞。

% !TEX root = ../eigenfaces_algorithm.tex
%=====================================================
\section{拼装：从 280 维的 $\mathbf{u}$ 到 4096 维的特征脸}

每个 Gram 侧特征对 $(\lambda_i,\mathbf{u}_i)$ 按第 6 节公式提升：
\[
\mathbf{v}_i=\frac{\mat{X}\mT\mathbf{u}_i}{\sqrt{\lambda_i}}\in\R^{4096},\qquad
i=1,\dots,k .
\]
三件实务：
\begin{enumerate}
  \item \textbf{丢掉伪根}：$\lambda_i<10^{-12}\lambda_1$ 的直接丢弃（数值噪声冒充的特征值，除以 $\sqrt\lambda$ 会爆炸）。
  \item \textbf{符号固定}：$\pm\mathbf{v}_i$ 都是合法特征向量。约定“绝对值最大的分量取正”，
  保证每次运行、每次与基准比对时方向一致，图像不会黑白颠倒。
  \item \textbf{排序}：按 $\lambda$ 降序拼成 $\mat{V}_k=[\mathbf{v}_1,\dots,\mathbf{v}_k]\in\R^{4096\times k}$，
  有 $\mat{V}_k\mT\mat{V}_k=\mat{I}_k$（正交性在 Gram 侧已保证，第 6.3 节验证过）。
\end{enumerate}

\textbf{与 \texttt{eigh} 比对的正确姿势。}基准真值不要用 4096² 的 \texttt{eigh}（不可行），
用 $\mat{X}$ 的瘦 SVD：$\lambda_i^{\text{ref}}=\sigma_i^2$、$\mathbf{v}_i^{\text{ref}}$ 取右奇异向量。三项指标：
\begin{enumerate}
  \item 特征值相对误差：$|\lambda_i-\lambda_i^{\text{ref}}|/\lambda_i^{\text{ref}}$（对 $i$ 画曲线，看尾部抬升）；
  \item 特征向量误差（先对齐符号）：$\min\big(\norm{\mathbf{v}_i-\mathbf{v}_i^{\text{ref}}},\ \norm{\mathbf{v}_i+\mathbf{v}_i^{\text{ref}}}\big)$；
  \item \textbf{子空间主角}：$\mat{V}_k\mT\mat{V}_k^{\text{ref}}$ 的奇异值即 $\cos\theta_i$。
  当两个特征值几乎重合时，单条特征向量会在子空间内部自由旋转、无法逐条比对，
  只有子空间角度是稳定的比较对象。
\end{enumerate}

% !TEX root = ../eigenfaces_algorithm.tex
%=====================================================
\section{投影与重建：一张脸 = 均值脸 + 若干份特征脸}

\subsection{正变换：提取坐标}

对任何一张脸（训练或测试都一样）：
\[
\mathbf{x}=\mathbf{a}-\bar{\mathbf{a}}\quad(\text{中心化})
\qquad\Longrightarrow\qquad
c_i=\inner{\mathbf{v}_i}{\mathbf{x}}=\mathbf{v}_i\mT\mathbf{x},\qquad
\mathbf{c}=\mat{V}_k\mT\mathbf{x}\in\R^{k}.
\]
$c_i$ 的含义：这张脸里含有多少份“第 $i$ 号变化模式”。
因为 $\{\mathbf{v}_i\}$ 是正交基，坐标唯一、且求法就是内积——
这正是课程里“基与坐标唯一性、基变换”的直接应用。

\subsection{逆变换：重建}

\[
\hat{\mathbf{x}}=\mat{V}_k\mathbf{c}
=\sum_{i=1}^k c_i\mathbf{v}_i
=\mat{V}_k\mat{V}_k\mT\,\mathbf{x}
\qquad\Longrightarrow\qquad
\boxed{\ \hat{\mathbf{a}}=\mat{V}_k\mat{V}_k\mT(\mathbf{a}-\bar{\mathbf{a}})+\bar{\mathbf{a}}\ }
\]

\begin{intuition}
重建公式读成一句话：\textbf{这张脸 $\approx$ 平均脸 + $c_1$ 份 1 号特征脸 + $c_2$ 份 2 号特征脸 + ……}。
基正交归一带来一个大便利：$\mat{V}_k^{-1}=\mat{V}_k\mT$，
从坐标回到图像不需要解任何方程组，转置乘回去即可（谱定理的红利）。
$k$ 越大，加进去的“模式”越多，脸越清晰；$k=279$ 时可完美复原（薄板被完整描述）。
\end{intuition}

\subsection{投影矩阵 $P=\mat{V}_k\mat{V}_k\mT$：它是一个真正的“投影”}

$\mat{P}=\mat{V}_k\mat{V}_k\mT$（作用在 $\R^{4096}$ 上）满足：
\[
\mat{P}^2=\mat{P}\qquad(\text{投影两次不再改变}),\qquad
\mat{P}\mT=\mat{P}\qquad(\text{对称}),\qquad
\operatorname{rank}\mat{P}=k=\operatorname{trace}\mat{P}.
\]
验证 $\mat{P}^2=\mat{P}$ 只需 $\mat{V}\mT\mat{V}=\mat{I}$：
$\mat{P}^2=\mat{V}(\mat{V}\mT\mat{V})\mat{V}\mT=\mat{V}\mat{V}\mT=\mat{P}$。
这三条性质对应课程中“正交投影 $P^2=P=P\mT$”的完整刻画。
残差
\[
\mathbf{r}=\mathbf{x}-\mat{P}\mathbf{x}
\]
满足 $\mathbf{v}_i\mT\mathbf{r}=c_i-c_i=0$：\textbf{残差与特征脸子空间垂直}，
即重建已经拿走了子空间里的一切，剩下的是子空间“管不着”的部分。

\subsection{误差有多大：毕达哥拉斯给出解析答案}

正交性使勾股定理成立：
\[
\norm{\mathbf{x}}^2=\norm{\mathbf{c}}^2+\norm{\mathbf{r}}^2
\qquad\text{其中}\quad
\norm{\mathbf{c}}^2=\sum_{i\le k}c_i^2 .
\]
对整个训练集求和（利用 $\sum_j\norm{\mathbf{x}_j}^2=\operatorname{trace}(\mat{X}\mT\mat{X})=\sum_i\lambda_i$，
以及 $\sum_j\sum_{i\le k}c_{ij}^2=\sum_{i\le k}\lambda_i$）：
\[
\text{平均每张图的残差平方和}
=\frac1n\sum_{j}\norm{\mathbf{r}_j}^2
=\frac{1}{n}\sum_{i>k}\lambda_i .
\qquad(\text{每像素 MSE 再除以 }d=4096)
\]
\textbf{重建误差 = 被扔掉的那些方向的方差之和。}
这条解析曲线应与实验实测 MSE 曲线重合；不重合说明特征向量或中心化有错——
它比“重建图看起来挺像”严格得多，建议作为报告的硬校验。

另一个引用级结论（Eckart--Young 定理，作为引用依据使用）：
在一切秩 $\le k$ 的逼近中，向特征脸子空间的正交投影使
$\norm{\mat{X}-\mat{X}_k}_F$ 最小，且最小值的平方恰为 $\sum_{i>k}\lambda_i$。
这就是“PCA 是最优线性降维”的准确含义。

\subsection{累计解释方差}

\[
R(k)=\frac{\sum_{i\le k}\lambda_i}{\sum_{i\le 279}\lambda_i}
\qquad\text{达到 }90\%\text{ 或 }95\%\text{ 的 }k\text{ 用插值从曲线上定位。}
\]
它回答的问题：\textbf{要保留九成五的“差异信息”，最少需要几张特征脸？}
人脸数据的典型答案是几十张——这正是第 1 节直觉“人的差异活在少数维度里”的定量版本。

% !TEX root = ../eigenfaces_algorithm.tex
%=====================================================
\section{识别：在 $k$ 维坐标里找最近邻}

\subsection{流程}

\begin{enumerate}
  \item \textbf{建画廊}：280 张训练图各自投影，得到坐标矩阵 $\mat{C}\in\R^{280\times k}$（每行一个 $\mathbf{c}_j$）及标签 $y_j$（人的编号）。
  \item \textbf{投测试图}：$\mathbf{c}_t=\mat{V}_k\mT(\mathbf{a}_t-\bar{\mathbf{a}})$。
  \item \textbf{找最近邻}：算 280 个距离 $d_j=\norm{\mathbf{c}_t-\mathbf{c}_j}_2$，取最小者的标签作为预测（平票取索引最小者）。
  \item \textbf{算准确率}：120 张测试图判对的比例。对 $k\in\{5,10,20,50,100\}$ 制表并画曲线。
\end{enumerate}

\begin{intuition}
为什么坐标距离能代替“长相相似度”？因为 $\mat{V}_k$ 的列是正交基，
内积、距离这些几何量在基变换下保持不变（正交变换不改变长度与夹角）。
$\norm{\mathbf{c}_t-\mathbf{c}_j}$ 本来就是“两张脸在特征脸子空间里的真实距离”，
只是把 4096 维里无信息的维度统统删掉了。同一个人的三张测试脸，
减去均值后剩下的“差异模式”与训练脸高度相似，因此在坐标空间里自然聚成一团。
\end{intuition}

\subsection{可选变体：白化距离}

经典 Eigenfaces 文献常用“白化”后的距离：
\[
d_j^2=\sum_{i\le k}\frac{(c_{t,i}-c_{j,i})^2}{\lambda_i} ,
\]
即先除掉每个方向的天然幅度（1 号特征脸系数天然比 50 号大几个数量级，不除会让头部方向垄断距离）。
它通常改变识别结果，\textbf{只能作为并列变体报告，不能与主结果混为一谈}。

\subsection{为什么 $k$ 不是越大越好（一条值得写进报告的解释）}

正交分解给出精确的恒等式：
\[
\norm{\mathbf{a}_t-\mathbf{a}_j}^2
=\norm{\mathbf{c}_t-\mathbf{c}_j}^2+\norm{\mathbf{r}_t-\mathbf{r}_j}^2 ,
\]
而 $k$ 维坐标距离\textbf{恰好扔掉了残差项} $\norm{\mathbf{r}_t-\mathbf{r}_j}^2$。
$k$ 小：残差大、且携带真实的判别信息（人脸轮廓差异还躺在被扔掉的方向里），
丢掉它是欠拟合；$k$ 大：残差项变小，但噪声维度的坐标进入距离并开始主导排序，
是过拟合。所以准确率随 $k$ 一般先升后平（或缓降），最佳 $k$ 常落在解释方差 $90\sim95\%$ 附近。
注意这是实验观察，报告中应表述为“实验中观察到”，而非理论必然。

% !TEX root = ../eigenfaces_algorithm.tex
%=====================================================
\section{收敛分析怎么呈现（对应题目的可衡量产出）}

画一张散点图：横轴 $|\lambda_{i+1}/\lambda_i|$，纵轴第 $i$ 对的迭代次数（对数轴）。
理论预测
\[
k_i\approx\frac{\ln(1/\varepsilon)}{\ln(\lambda_i/\lambda_{i+1})}
\]
是一条确定的曲线，把它画上去与实测点对照：
\begin{itemize}
  \item \textbf{吻合段}：谱间隙主导，幂迭代行为与理论一致；
  \item \textbf{偏离段}：从某个 $i$ 开始实测迭代数高于理论线——这正是收缩误差累积开始压过谱间隙效应的位置，
  两张图（迭代数 vs 间隙、相对误差 vs $i$）互为印证。
\end{itemize}
这张图是本题目最有“研究味”的产出：它把\textbf{课程里的谱半径收敛理论}
（Jacobi 迭代 $\rho(B)<1$ 的类比）迁移到了特征值算法上。

% !TEX root = ../eigenfaces_algorithm.tex
%=====================================================
\section{交互演示与作图注意事项}

演示逻辑：预处理只做一次——算出 $\bar{\mathbf{a}}$、全部 279 个特征对
（前 50 个来自幂迭代 + 收缩，其余来自 \texttt{eigh}）、缓存全部坐标与各 $k$ 的准确率。
滑块只切 $k$，实时重算 $\mathbf{c}=\mat{V}_k\mT\mathbf{x}$、$\hat{\mathbf{a}}=\mat{V}_k\mathbf{c}+\bar{\mathbf{a}}$，
同屏显示：原图、重建图、前 16 张特征脸，并读出 $R(k)$ 与该 $k$ 的测试准确率。

\textbf{特征脸的显示规范（重要）：}特征向量的分量是实数、有正有负，
必须使用统一色标（灰度中点 $=0$，或按全体分量的全局分位裁剪）。
若每张图各自做 min--max 归一化，16 张图会全都呈现“差不多亮”的假象——
这是初学者报告里最常见的视觉误导。

% !TEX root = ../eigenfaces_algorithm.tex
%=====================================================
\section{验证清单：怎么证明算法没写错}

\begin{enumerate}
  \item \textbf{正交归一}：$\norm{\mat{V}_k\mT\mat{V}_k-\mat{I}_k}_F$ 应在 $10^{-10}$ 量级
  （\texttt{eigh} 基准约 $10^{-15}$）。漂移大说明再正交化没做到位。
  \item \textbf{投影性质}：取随机向量 $\mathbf{z}$，验证
  $\norm{\mat{P}(\mat{P}\mathbf{z})-\mat{P}\mathbf{z}}/\norm{\mat{P}\mathbf{z}}\approx 0$
  与 $|\mathbf{z}\mT\mat{P}\mathbf{z}-(\mat{P}\mathbf{z})\mT(\mat{P}\mathbf{z})|\approx 0$。
  用随机向量检验即可，\textbf{不要真的构造 $4096\times4096$ 的 $\mat{P}$}。
  \item \textbf{秩}：$\operatorname{trace}\mat{P}=\sum_i\norm{\mathbf{v}_i}^2$ 应精确等于 $k$。
  \item \textbf{特征残差}：$\norm{\mat{G}\mathbf{u}_i-\lambda_i\mathbf{u}_i}$ 与
  $\norm{\mat{C}\mathbf{v}_i-\frac{\lambda_i}{n-1}\mathbf{v}_i}$。
  \item \textbf{满秩重建}：$k=279$ 时 $\norm{\mat{X}-\mat{V}_k\mat{V}_k\mT\mat{X}}_F/\norm{\mat{X}}_F$
  应达机器精度量级——这是“Gram $\to\mathbf{u}\to\mathbf{v}\to$ 投影”整条链路无误的最强证据。
  \item \textbf{MSE 对解析曲线}：实测重建 MSE 应与 $\sum_{i>k}\lambda_i/(n\,d)$ 重合（第 9.4 节）。
\end{enumerate}

% !TEX root = ../eigenfaces_algorithm.tex
%=====================================================
\section{常见坑清单}

\begin{center}
\begin{tabular}{p{0.44\textwidth}p{0.46\textwidth}}
\toprule
\textbf{错误做法} & \textbf{后果} \\
\midrule
用测试集算均值 / 中心化 & 信息泄漏，测试坐标系与训练基不一致 \\
重建后忘记加回 $\bar{\mathbf{a}}$ & 重建图整体发暗、无法与原图比较 \\
特征向量忘记除以 $\sqrt{\lambda}$ & 非单位基，坐标整体缩放，1-NN 距离被扭曲 \\
收缩时不做再正交化 & 第 5 个之后全部“退化回第一主方向” \\
\texttt{eigh} 升序 vs 幂迭代降序直接比对 & 全部错位 \\
比对特征向量不处理 $\pm$ 符号 & 误差恒为 $\sqrt{2}$，无法判断好坏 \\
只用 $|\Delta\hat\lambda|$ 判停 & 方向未收敛就退出（Rayleigh 商二阶乐观）\\
把 $\lambda$ 当成方差 & 绝对值差 $n-1=279$ 倍（比值不受影响）\\
像素 $0\sim255$ 与 $0\sim1$ 混用 & MSE 差 $65025$ 倍 \\
$k>279$ & $\lambda=0$，除以 $\sqrt{\lambda}$ 得 NaN \\
特征脸逐图各自归一化显示 & 视觉上全部“一样亮”，误导解读 \\
\bottomrule
\end{tabular}
\end{center}

% !TEX root = ../eigenfaces_algorithm.tex
%=====================================================
\section*{附录 A：符号速查表}

\begin{center}
\begin{tabular}{ll}
\toprule
符号 & 含义 \\
\midrule
$n=280,\ d=4096$ & 训练脸数、像素维数 \\
$\mat{A}\in\R^{n\times d}$ & 原始数据矩阵（每行一张脸）\\
$\bar{\mathbf{a}}$ & 均值脸（只由训练集计算）\\
$\mat{X}=\mat{A}-\mathbf{1}\bar{\mathbf{a}}\mT$ & 中心化数据矩阵，$\operatorname{rank}\le 279$ \\
$\mat{G}=\mat{X}\mat{X}\mT\in\R^{n\times n}$ & Gram（相似度）矩阵 \\
$\lambda_i,\ \mathbf{u}_i$ & $\mat{G}$ 的特征值 / 特征向量（280 维）\\
$\mathbf{v}_i=\mat{X}\mT\mathbf{u}_i/\sqrt{\lambda_i}$ & 特征脸（4096 维）\\
$\mat{V}_k=[\mathbf{v}_1\cdots\mathbf{v}_k]$ & 前 $k$ 张特征脸，$\mat{V}_k\mT\mat{V}_k=\mat{I}$ \\
$\mathbf{c}=\mat{V}_k\mT\mathbf{x}$ & 一张脸在特征脸基下的 $k$ 维坐标 \\
$\mat{P}=\mat{V}_k\mat{V}_k\mT$ & 正交投影矩阵，$\mat{P}^2=\mat{P}=\mat{P}\mT$ \\
$\hat\lambda=\mathbf{u}\mT\mat{G}\mathbf{u}$ & Rayleigh 商（特征值的二阶精确估计）\\
$R(k)$ & 累计解释方差比 \\
\bottomrule
\end{tabular}
\end{center}

% !TEX root = ../eigenfaces_algorithm.tex
%=====================================================
\section*{附录 B：整体流程一览}

\begin{enumerate}
  \item 载入 400 张脸，分层划分 280/120；像素归一化到 $[0,1]$。
  \item 训练集求 $\bar{\mathbf{a}}$，得中心化 $\mat{X}\in\R^{280\times4096}$。
  \item 计算 $\mat{G}=\mat{X}\mat{X}\mT$（或每步以 $\mathbf{u}\mapsto\mat{X}(\mat{X}\mT\mathbf{u})$ 隐式计算）。
  \item 幂迭代 + 再正交化 + Hotelling 收缩，取前 $k\le 50$ 对 $(\lambda_i,\mathbf{u}_i)$；$k>50$ 用 \texttt{eigh}。
  \item 提升 $\mathbf{v}_i=\mat{X}\mT\mathbf{u}_i/\sqrt{\lambda_i}$，固定符号、降序拼装 $\mat{V}_k$。
  \item 验证：正交性、$\mathtt{trace}\,\mathtt{P}=k$、投影幂等、特征残差、$k=279$ 满秩重建、与 SVD 基准比对。
  \item 投影/重建：$\mathbf{c}=\mat{V}_k\mT(\mathbf{a}-\bar{\mathbf{a}})$，$\hat{\mathbf{a}}=\mat{V}_k\mathbf{c}+\bar{\mathbf{a}}$；
        画 MSE--$k$ 与 $R(k)$ 曲线并与解析值 $\sum_{i>k}\lambda_i$ 对照。
  \item 1-NN 识别：$k$ 维坐标欧氏距离，对 $k\in\{5,10,20,50,100\}$ 报准确率。
  \item 收敛分析：迭代次数 vs $|\lambda_{i+1}/\lambda_i|$ 散点 + 理论线对照。
  \item 交互演示：滑块切 $k$，展示原图 / 重建 / 前 16 张特征脸（统一色标）。
\end{enumerate}


\end{document}
